\documentclass[12pt]{article}
\usepackage[utf8]{inputenc}
\usepackage[T1]{fontenc}
\usepackage[margin=1in]{geometry}
\usepackage{setspace}
\setlength{\parindent}{2em}
\setlength{\parskip}{1.25em}
\renewcommand{\baselinestretch}{1.0}
\makeatletter
\newcommand{\@subtitle}{}
\newcommand{\subtitle}[1]{\renewcommand{\@subtitle}{#1}}
\renewcommand{\@maketitle}{%
  \newpage
  \null
  \vskip 1em%
  \begin{center}%
    {\LARGE \@title \par}%
    \ifx\@subtitle\@empty\else
      \vskip 0.5em%
      {\large \@subtitle \par}%
    \fi
  \end{center}%
  \par
  \vskip 1em}
\makeatother
\usepackage{xcolor}
\usepackage[colorlinks=true,urlcolor=blue!55!black,linkcolor=black]{hyperref}

\title{\textit{Barron v.\ Mayor \& City Council of Baltimore}, 32 U.S.\ (7 Pet.) 243 (1833)}
\subtitle{Case Brief}
\author{}
\date{}

\begin{document}
\maketitle

\section{Memory Jogger}

In \href{https://supreme.justia.com/cases/federal/us/32/243/}{\textit{Barron v. Baltimore}}, Chief Justice John Marshall held for a unanimous Court that the Fifth Amendment's Takings Clause---and by extension the entire federal Bill of Rights---restrains only the federal government and does not apply to state or municipal governments.

\section{Detailed Case Facts}

\textbf{Factual Background.} John Barron and his business partner John Craig (of whom Barron was the surviving partner) owned an extensive and highly profitable commercial wharf in the eastern section of Baltimore along the Patapsco River. At the time they purchased the wharf, it enjoyed the deepest water in Baltimore harbor, enabling large commercial vessels to dock there. Between 1815 and 1821, the Mayor and City Council of Baltimore enacted a series of municipal ordinances to pave and regrade city streets, regulate the harbor, and protect public health. In carrying out these public works---through new street grades, paving, mounds, embankments, and other artificial structures---the city diverted several natural streams flowing from the hills bordering Baltimore so that they emptied directly toward Barron's wharf. During heavy rains, the swollen streams washed large quantities of sand and earth down from the hills and deposited them along and in front of the wharf. The accumulated silt made the water so shallow that vessels of substantial burden could no longer reach the wharf, destroying its commercial income and rendering it nearly worthless (\href{https://supreme.justia.com/cases/federal/us/32/243/}{32 U.S. at 243--44}).

\textbf{Statutory and Constitutional Provisions.} The city never offered or paid compensation to Barron. Instead, it justified its street-grading and drainage projects under the Baltimore city charter (Maryland Acts of 1796, ch. 68) and subsequent Maryland legislative statutes delegating municipal authority over streets, harbor waters, and city health. Barron contended that the city's state-authorized destruction of his wharfage rights violated the Takings Clause of the Fifth Amendment to the U.S. Constitution, which declares that ``private property [shall not] be taken for public use, without just compensation.'' Because Barron sought Supreme Court review of a state court judgment, federal appellate jurisdiction depended on \textbf{Section 25 of the Judiciary Act of 1789}, which authorized the Supreme Court to review final state-court decisions upholding state statutes or state-delegated authority against a claim of repugnancy to the Federal Constitution.

\textbf{Historical Context.} Decided during the final years of Chief Justice Marshall's tenure and amid the Nullification Crisis of 1832--1833, \textit{Barron} squarely tested the antebellum structure of dual federalism. During the 1787--1788 ratification debates, Anti-Federalists had feared an overbearing national government and demanded a federal bill of rights; in 1789, James Madison had proposed an amendment expressly prohibiting States from infringing conscience, speech, press, or jury trial (which Madison called ``the most valuable of the whole list''), but the Senate rejected that state-binding provision (\href{https://www.law.cornell.edu/constitution-conan/amendment-14/overview-of-incorporation-of-the-bill-of-rights}{Cornell LII Constitution Annotated}). When Barron sued in 1822, the Fourteenth Amendment did not yet exist, and most individual liberties against state abuse were left to state constitutions and state courts.

\section{Procedural History}

\textbf{Trial Court.} In 1822, John Barron filed an action on the case in the Baltimore County Court against the Mayor and City Council of Baltimore seeking damages for the destruction of his wharf. The city defended on two principal grounds: (1) its actions were lawfully authorized by Maryland statutes and its municipal charter, and (2) Barron's complaint alleged at most a non-actionable public nuisance rather than a special individual injury. The Baltimore County Court ruled against the city on all points, and a jury awarded Barron a verdict of \$4,500 (\href{https://supreme.justia.com/cases/federal/us/32/243/}{32 U.S. at 244}).

\textbf{State Appellate Court.} The city appealed to the Court of Appeals for the Western Shore of the State of Maryland (then Maryland's highest court), which reversed the Baltimore County Court's judgment and refused to remand the case for a new trial.

\textbf{U.S. Supreme Court and Disposition.} Barron prosecuted a writ of error to the U.S. Supreme Court under Section 25 of the Judiciary Act of 1789. The case was argued on February 8 and 11, 1833. On February 16, 1833, the Supreme Court held that because the Fifth Amendment does not apply to state legislation or municipal actions authorized thereby, no federal constitutional question was presented. The Court therefore \textbf{dismissed the writ of error for want of jurisdiction} (\href{https://supreme.justia.com/cases/federal/us/32/243/}{32 U.S. at 251}).

\section{Judicial Votes}

The Court decided the case \textbf{unanimously (7--0)} (\href{https://supreme.justia.com/cases/federal/us/32/243/}{32 U.S. at 247}):

\begin{itemize}
  \item \textbf{Majority (7):} Chief Justice John Marshall delivered the opinion of the Court, joined by Associate Justices William Johnson, Gabriel Duvall, Joseph Story, Smith Thompson, John McLean, and Henry Baldwin.
  \item \textbf{Concurrences:} None.
  \item \textbf{Dissents:} None.
\end{itemize}

\section{Holding}

\textbf{Judgment.} The writ of error is dismissed for want of jurisdiction because there is no repugnancy between the Maryland statutes relied upon by Baltimore and the Constitution of the United States.

\textbf{Rule Announced.} The Takings Clause of the Fifth Amendment---and, by the opinion's reasoning, the provisions of the Bill of Rights generally---operates solely as a limitation on the powers of the federal government and does not bind state or local governments. Consequently, constitutional limitations expressed in general terms apply only to the federal government unless the text expressly directs them to the States.

\textbf{Holding vs. Dicta.} Strictly speaking, the Court's \textbf{holding} was confined to the Fifth Amendment's Just Compensation Clause and its jurisdictional consequence under Section 25 of the Judiciary Act of 1789. However, Chief Justice Marshall's broader structural and historical rationale---declaring that ``[t]hese amendments contain no expression indicating an intention to apply them to the State governments''---functioned as sweeping \textbf{dicta} as to the other amendments in the Bill of Rights, which the Court swiftly treated as settled law across the entire Bill of Rights in subsequent antebellum decisions (see \href{https://www.law.cornell.edu/constitution-conan/amendment-14/overview-of-incorporation-of-the-bill-of-rights}{Cornell LII Constitution Annotated}, citing \textit{Permoli v. Municipality No. 1}, 44 U.S. (3 How.) 589 (1845), and \textit{Twitchell v. Commonwealth}, 74 U.S. (7 Wall.) 321 (1869)). While \textit{Barron} has never been formally overruled as an interpretation of the original 1791 Bill of Rights standing alone, its practical effect was largely superseded after 1868 by the selective incorporation of nearly all Bill of Rights guarantees against the States through the Fourteenth Amendment's Due Process Clause, beginning with the Takings Clause itself in \href{https://supreme.justia.com/cases/federal/us/166/226/}{\textit{Chicago, Burlington \& Quincy Railroad Co. v. Chicago}, 166 U.S. 226 (1897)}.

\section{Analysis of Opinions}

\subsection{Opinion of the Court (Chief Justice Marshall)}

Chief Justice Marshall described the issue as ``of great importance, but not of much difficulty'' and organized his unanimous opinion around four tightly linked arguments: structural theory, intratextual comparison, institutional logic, and ratification history (\href{https://supreme.justia.com/cases/federal/us/32/243/}{32 U.S. at 247--51}).

\textbf{1. Social-Compact and Structural Theory.} Marshall began from first principles of federalism. The Federal Constitution was ordained by the people of the United States ``for their own government, and not for the government of the individual States.'' Each State already possessed its own constitution containing whatever limits on state power its citizens deemed appropriate. Therefore, limitations on power stated in general terms in the Federal Constitution are ``naturally, and we think necessarily, applicable to the government created by the instrument,'' not to distinct state governments created by different instruments.

\textbf{2. Intratextual Comparison of Article I, \S\S~9 and 10.} Barron's counsel argued that Article I, \S~10 proved the original Constitution restrained States to protect citizens' liberties. Marshall turned that argument against Barron. Article I, \S~9 uses general, passive language---such as ``No Bill of Attainder or ex post facto Law shall be passed''---without naming Congress, yet unquestionably binds only the federal government. We know this conclusively because the very next section, Article I, \S~10, repeats the exact prohibition using express state-binding language: ``No State shall \dots pass any Bill of Attainder, ex post facto Law.'' Whenever the Framers intended a constitutional restriction to bind the States (typically on matters of national concern such as treaties, coinage, or war), they said so ``in positive words.'' Because the Fifth Amendment omits any phrase like ``No State shall,'' it cannot be read to depart from that drafting convention.

\textbf{3. Institutional Logic of Article V.} Marshall bolstered his textual reading with practical common sense. If the citizens of a particular State wanted greater constitutional protection against their own state legislature, they could simply amend their own state constitution via a local convention. ``The unwieldy and cumbrous machinery'' of Article V---requiring two-thirds of Congress and three-fourths of the States---``could never have occurred to any human being as a mode of doing that which might be effected by the State itself.''

\textbf{4. Ratification History.} Finally, Marshall appealed to living memory (he had personally participated in Virginia's 1788 ratifying convention). The Bill of Rights was proposed and ratified to quiet widespread Anti-Federalist fears that the newly created national government might abuse its powers---``not against those of the local governments.'' Because the amendments contained no expression indicating an intent to bind the States, the Court could not extend them by judicial construction. Finding no conflict between Maryland's statutes and the Federal Constitution, Marshall dismissed the writ of error for lack of jurisdiction.

\section{Examples: Future Applications}

\begin{enumerate}
  \item \textbf{Same-Side (Direct Application of \textit{Barron} --- Unincorporated Bill of Rights Provision):} A State prosecutes a defendant for first-degree murder by a prosecutor's ``information'' following a magistrate's hearing, rather than by a grand jury indictment as required by the Fifth Amendment. \textit{Outcome:} Permitted. Under \textit{Barron}, the Fifth Amendment Grand Jury Clause does not directly bind the States, and in \href{https://supreme.justia.com/cases/federal/us/110/516/}{\textit{Hurtado v. California}, 110 U.S. 516 (1884)}, the Supreme Court declined to incorporate it through the Fourteenth Amendment Due Process Clause---one of the few Bill of Rights guarantees where \textit{Barron}'s rule still controls state practice today.

  \item \textbf{Same-Side (Direct Application of \textit{Barron} --- Civil Jury Trials in State Court):} A State abolishes civil jury trials in state-court tort actions exceeding twenty dollars and directs that all such claims be tried to a judge. A plaintiff invokes the Seventh Amendment's Civil Jury Trial Clause. \textit{Outcome:} Permitted under the Federal Constitution. Because the Seventh Amendment does not directly apply to the States under \textit{Barron} and has never been incorporated through the Fourteenth Amendment, States remain free to structure civil trial procedures under their own constitutions.

  \item \textbf{Opposite-Side (Superseded via Fourteenth Amendment Incorporation --- State Taking Without Compensation):} A municipality condemns a railroad's right-of-way to open a public street and awards nominal compensation of \$1, and the landowner challenges the state judgment under the Fourteenth Amendment. \textit{Outcome:} Federal constitutional protection applies. While \textit{Barron} barred direct reliance on the Fifth Amendment, the Supreme Court held in \href{https://supreme.justia.com/cases/federal/us/166/226/}{\textit{Chicago, Burlington \& Quincy Railroad Co. v. Chicago}, 166 U.S. 226 (1897)}, that a State's taking of private property for public use without just compensation violates the Due Process Clause of the Fourteenth Amendment (see \href{https://www.law.cornell.edu/constitution-conan/amendment-5/overview-of-the-takings-clause}{Cornell LII Takings Clause Overview}).

  \item \textbf{Opposite-Side (Superseded via Selective Incorporation --- Modern Bill of Rights Claims Against States):} A city bans handgun possession in the home, or a State seizes a \$42,000 vehicle as civil forfeiture for a minor drug offense carrying a maximum \$10,000 fine. \textit{Outcome:} Federal constitutional limits bind the State and city. Even though \textit{Barron} held the original Bill of Rights inapplicable to States, the Supreme Court incorporated the Second Amendment in \href{https://supreme.justia.com/cases/federal/us/561/742/}{\textit{McDonald v. City of Chicago}, 561 U.S. 742 (2010)}, and the Eighth Amendment Excessive Fines Clause in \href{https://supreme.justia.com/cases/federal/us/586/17-1091/}{\textit{Timbs v. Indiana}, 586 U.S. 146 (2019)}, via the Fourteenth Amendment.

  \item \textbf{Fence-Sitter (Borderline --- Unresolved Incorporation of the Third Amendment):} During a protracted civil emergency, a State governor orders National Guard troops operating under state active duty (not federalized) to quarter inside private residences without the owners' consent. \textit{Why It Is Close:} The Supreme Court has never squarely decided whether the Third Amendment is incorporated against the States (though the Second Circuit held it was in \textit{Engblom v. Carey}, 677 F.2d 957 (2d Cir. 1982)). Under \textit{Barron} alone, the Third Amendment would not bind state troops; under modern \textit{McDonald}/\textit{Timbs} incorporation doctrine, a court would ask whether freedom from nonconsensual quartering is ``deeply rooted in this Nation's history and tradition'' and ``fundamental to our scheme of ordered liberty.''
\end{enumerate}

\section{Critique}

\textbf{Textualist and Originalist Perspectives.} Most modern originalists praise \textit{Barron} as an exemplar of disciplined intratextualism and historical fidelity. Comparing the general language of Article I, \S~9 with the express ``No State shall'' language of Article I, \S~10---and corroborating that reading with the First Congress's rejection of Madison's proposed state-binding amendment---strongly supports Marshall's conclusion that the 1791 Bill of Rights bound only the federal government. Yet \textit{Barron} is not immune from textualist critique. Professor William Winslow Crosskey famously attacked \textit{Barron} as a sharp departure from the plain text, pointing out that the First Amendment explicitly names ``Congress'' (``Congress shall make no law'') whereas the Second through Eighth Amendments use unrestricted passive language (``shall not be infringed,'' ``shall not be violated''), implying by \textit{expressio unius} that the general amendments bound both federal and state governments---a thesis vigorously debated by Professor Charles Fairman in \href{https://chicagounbound.uchicago.edu/cgi/viewcontent.cgi?article=2869&context=uclrev}{\textit{A Reply to Professor Crosskey}}.

More subtly, modern originalists such as Akhil Reed Amar (\href{https://scholarship.law.wm.edu/cgi/viewcontent.cgi?article=1398&context=wmborj}{discussed in William \& Mary Bill of Rights Journal}) and Kurt Lash emphasize that \textit{Barron} itself shaped the drafting of the Fourteenth Amendment: Representative John Bingham expressly cited \textit{Barron} in 1866 to explain why a new constitutional amendment containing explicit ``No State shall'' language was necessary to make fundamental rights enforceable against the States through the Privileges or Immunities Clause. From this perspective, the real historical error was not \textit{Barron}, but the Court's evisceration of the Privileges or Immunities Clause in the \textit{Slaughter-House Cases}, 83 U.S. 36 (1873), which forced the Court to channel incorporation awkwardly through ``substantive'' Due Process (as Justice Thomas argued in his \href{https://supreme.justia.com/cases/federal/us/561/742/}{\textit{McDonald}} concurrence).

\textbf{Progressive and Rights-Protective Perspectives.} Progressive critics highlight the stark human and political costs of \textit{Barron}'s dual-federalism premise. Marshall's cheerful assurance that citizens dissatisfied with state abuses could simply convene a state constitutional convention ignored the reality of antebellum America: enslaved people, free Black citizens, women, religious dissenters, and abolitionists lacked the franchise or political power to reform oppressive state governments. By immunizing States from the First and Fifth Amendments, \textit{Barron} left Southern States free to criminalize abolitionist speech and press before the Civil War and delayed national enforcement of civil liberties until the twentieth-century incorporation revolution inaugurated in \textit{Gitlow v. New York}, 268 U.S. 652 (1925) (\href{https://www.law.cornell.edu/constitution-conan/amendment-14/overview-of-incorporation-of-the-bill-of-rights}{Cornell LII Constitution Annotated}).

\textbf{Logical Weaknesses.} Two analytical tensions stand out in Marshall's opinion. First, Marshall's claim that Article I, \S~10 merely restrained States on matters of national concern (like treaties and coinage) glossed over \S~10's Contract Clause, Bill of Attainder Clause, and Ex Post Facto Clause---provisions designed squarely to protect individual property and liberty from state abuse, as Marshall himself had forcefully enforced in \textit{Fletcher v. Peck} (1810) and \textit{Dartmouth College v. Woodward} (1819). Second, Marshall treated the Bill of Rights as purely a structural bargain between sovereigns rather than a declaration of natural, pre-existing rights binding all American governments---a view later challenged by antebellum ``\textit{Barron} contrarians'' and ultimately transformed by Reconstruction.

\section{Key Quotations}

\begin{quote}
``The question thus presented is, we think, of great importance, but not of much difficulty. The Constitution was ordained and established by the people of the United States for themselves, for their own government, and not for the government of the individual States. Each State established a constitution for itself, and in that constitution provided such limitations and restrictions on the powers of its particular government as its judgment dictated.''\\
--- Chief Justice Marshall, Opinion of the Court, \href{https://supreme.justia.com/cases/federal/us/32/243/}{32 U.S. at 247}.
\end{quote}

\begin{quote}
``The powers they conferred on this government were to be exercised by itself, and the limitations on power, if expressed in general terms, are naturally, and we think necessarily, applicable to the government created by the instrument. They are limitations of power granted in the instrument itself, not of distinct governments framed by different persons and for different purposes.''\\
--- Chief Justice Marshall, Opinion of the Court, \href{https://supreme.justia.com/cases/federal/us/32/243/}{32 U.S. at 247}.
\end{quote}

\begin{quote}
``If the original Constitution, in the ninth and tenth sections of the first article, draws this plain and marked line of discrimination between the limitations it imposes on the powers of the General Government and on those of the State; if, in every inhibition intended to act on State power, words are employed which directly express that intent; some strong reason must be assigned for departing from this safe and judicious course in framing the amendments before that departure can be assumed. We search in vain for that reason.''\\
--- Chief Justice Marshall, Opinion of the Court, \href{https://supreme.justia.com/cases/federal/us/32/243/}{32 U.S. at 249}.
\end{quote}

\begin{quote}
``These amendments demanded security against the apprehended encroachments of the General Government --- not against those of the local governments. In compliance with a sentiment thus generally expressed, to quiet fears thus extensively entertained, amendments were proposed by the required majority in Congress and adopted by the States. These amendments contain no expression indicating an intention to apply them to the State governments. This court cannot so apply them.''\\
--- Chief Justice Marshall, Opinion of the Court, \href{https://supreme.justia.com/cases/federal/us/32/243/}{32 U.S. at 250}.
\end{quote}

\begin{quote}
``We are of opinion that the provision in the Fifth Amendment to the Constitution declaring that private property shall not be taken for public use without just compensation is intended solely as a limitation on the exercise of power by the Government of the United States, and is not applicable to the legislation of the States.''\\
--- Chief Justice Marshall, Opinion of the Court, \href{https://supreme.justia.com/cases/federal/us/32/243/}{32 U.S. at 250--51}.
\end{quote}

\end{document}
